\section{\texorpdfstring{هفته‌ی 11: روش‌های \lr{Monte Carlo} و الگوریتم
\lr{Metropolis\lr{--}Hastings}}{هفته‌ی 11: روش‌های Monte Carlo و الگوریتم Metropolis\lr{--}Hastings}}

\begin{quote}
\textbf{فصل 5: برآورد پارامتر با روش‌های داده‌محور \lr{(MCMC)}} · درس 1
ساعت · \lr{lab} 2 ساعت \textbf{داده‌ی اصلی:} داده‌های فراموشی هفته‌ی 2 با
پیشین‌های هفته‌ی 10، تا بتوانیم نمونه‌گیر را با پسین شبکه‌ای‌ای که از قبل
داریم مقایسه کنیم.
\end{quote}

\begin{center}\rule{0.5\linewidth}{0.5pt}\end{center}

\subsection{واژه‌های کلیدی}

واژه‌های مهم این درس با معادل فارسی‌شان (حتی اگر در درس‌های قبل آمده
باشند).

\begin{longtable}[]{@{}
  >{\raggedright\arraybackslash}p{(\linewidth - 4\tabcolsep) * \real{0.3333}}
  >{\raggedright\arraybackslash}p{(\linewidth - 4\tabcolsep) * \real{0.3333}}
  >{\raggedright\arraybackslash}p{(\linewidth - 4\tabcolsep) * \real{0.3333}}@{}}
\caption{\centering واژه‌های کلیدی}\tabularnewline
\toprule\noalign{}
\begin{minipage}[b]{\linewidth}\raggedright
\lr{English}
\end{minipage} & \begin{minipage}[b]{\linewidth}\raggedright
فارسی
\end{minipage} & \begin{minipage}[b]{\linewidth}\raggedright
توضیح کوتاه
\end{minipage} \\
\midrule\noalign{}
\endfirsthead
\toprule\noalign{}
\begin{minipage}[b]{\linewidth}\raggedright
\lr{English}
\end{minipage} & \begin{minipage}[b]{\linewidth}\raggedright
فارسی
\end{minipage} & \begin{minipage}[b]{\linewidth}\raggedright
توضیح کوتاه
\end{minipage} \\
\midrule\noalign{}
\endhead
\bottomrule\noalign{}
\endlastfoot
\lr{Monte Carlo} & \lr{Monte Carlo} (مونت‌کارلو) & تقریب با میانگین
نمونه‌های تصادفی \\
\lr{Metropolis algorithm} & الگوریتم \lr{Metropolis} & پیشنهاد یک حرکت و
پذیرش یا رد آن \\
\lr{parameter} & پارامتر & کمیت قابل تنظیم مدل که بین افراد یا شرایط فرق
می‌کند \\
\lr{estimate} / \lr{estimator} & برآورد / برآوردگر & مقدار به‌دست‌آمده
برای پارامتر / روش به دست آوردن آن \\
\lr{Markov chain Monte Carlo (MCMC)} & زنجیره‌ی مارکوف \lr{Monte Carlo} &
نمونه‌گیری از پسین با یک زنجیره‌ی تصادفی \\
\lr{forgetting curve} & منحنی فراموشی & کاهش یادآوری با گذشت زمان \\
\lr{prior} & توزیع پیشین \lr{(prior)} & باور درباره‌ی پارامتر پیش از دیدن
داده \\
\lr{sampler} & نمونه‌گیر \lr{(sampler)} & الگوریتمی که از توزیع نمونه
می‌گیرد \\
\lr{grid approximation} & تقریب شبکه‌ای \lr{(grid)} & محاسبه‌ی پسین روی
شبکه‌ای از مقادیر \\
\lr{posterior} & توزیع پسین \lr{(posterior)} & باور درباره‌ی پارامتر بعد
از دیدن داده \\
\lr{chain} & زنجیره & دنباله‌ی نمونه‌های \lr{MCMC} \\
\lr{proposal} & پیشنهاد \lr{(proposal)} & مقدار جدیدی که الگوریتم امتحان
می‌کند \\
\lr{acceptance rate} & نرخ پذیرش & نسبت پیشنهادهای پذیرفته‌شده \\
\lr{trace plot} & نمودار \lr{trace} & مقدار نمونه‌ها در برابر شماره‌ی
تکرار \\
\lr{autocorrelation} & خودهمبستگی & وابستگی نمونه‌های پشت سر هم \\
\lr{burn-in} / \lr{warm-up} & دوره‌ی گرم‌شدن \lr{(burn-in)} & نمونه‌های
اولیه که دور ریخته می‌شوند \\
\lr{sensitivity (d′)} & حساسیت & فاصله‌ی توزیع علامت و نویز \\
\lr{likelihood} & درست‌نمایی & احتمال داده به‌عنوان تابعی از پارامترها \\
\lr{starting point} & نقطه‌ی شروع & سوگیری اولیه به سمت یکی از پاسخ‌ها \\
\lr{response bias} & سوگیری پاسخ & تمایل به یک پاسخ، مستقل از محرک \\
\lr{large language model (LLM)} & مدل بزرگ زبانی & شبکه‌ی عصبی‌ای که کلمه‌ی
بعدی متن را پیش‌بینی می‌کند \\
\lr{self-consistency} & \lr{self-consistency} (رأی اکثریت) & چند بار
نمونه‌گیری جواب و انتخاب رأی اکثریت \\
\lr{simulation} & شبیه‌سازی & تولید داده‌ی مصنوعی با اجرای مدل \\
\lr{MCSE} & خطای استاندارد \lr{Monte Carlo} & خطای شبیه‌سازی در میانگین
پسین \\
\end{longtable}

\begin{center}\rule{0.5\linewidth}{0.5pt}\end{center}

\subsection{اهداف یادگیری}

در پایان این هفته، دانشجویان می‌توانند:

\begin{enumerate}
\def\labelenumi{\arabic{enumi}.}
\tightlist
\item
  \textbf{توضیح دهند} چرا کمیت‌های پسین انتگرال‌اند، و چطور نمونه‌گیری
  \lr{Monte Carlo} انتگرال‌ها را به میانگین تبدیل می‌کند، با خطایی که مثل
  \(1/\sqrt{S}\) کم می‌شود.
\item
  منطق \lr{Markov chain Monte Carlo} را \textbf{توصیف کنند}: زنجیره‌ای که
  توزیع بلندمدتش همان پسین است و فقط از پسین \emph{نرمال‌نشده} استفاده
  می‌کند.
\item
  الگوریتم \lr{Metropolis} (و \lr{Metropolis\lr{--}Hastings)} را برای یک مدل
  شناختی یک‌پارامتری و دوپارامتری از صفر \textbf{صورت‌بندی کنند}.
\item
  توزیع پیشنهادی \lr{(proposal)} را با استفاده از نرخ پذیرش، \lr{trace
  plot} و خودهمبستگی \textbf{تنظیم کنند} و دوره‌ی \lr{burn-in} را
  \textbf{مشخص کنند}.
\item
  خلاصه‌های پسین و کمیت‌های مشتق‌شده (مثلاً نیمه‌عمر) را از نمونه‌ها
  \textbf{حساب کنند} و نمونه‌گیر را با یک جواب معلوم \textbf{اعتبارسنجی
  کنند}.
\end{enumerate}

\begin{center}\rule{0.5\linewidth}{0.5pt}\end{center}

\subsection{مقدمه}

هر چیزی که از یک پسین می‌خواهیم یک \textbf{انتگرال} است: میانگین پسین
\(\int \theta\, p(\theta\mid\mathbf{y})\,d\theta\)، پسین حاشیه‌ای (با
انتگرال‌گیری روی بقیه‌ی پارامترها)، \(P(d' > 1\mid\mathbf{y})\)، توزیع
پیش‌بینی پسین، و \lr{marginal likelihood.} در هفته‌های 9\lr{--}10 این انتگرال‌ها
را روی شبکه حساب کردیم. این برای دو پارامتر کار کرد. یک مدل \lr{SDT}
سلسله‌مراتبی برای 20 شرکت‌کننده (هفته‌ی 13) بیش از 40 پارامتر دارد؛ شبکه‌ای
با فقط 10 نقطه برای هر پارامتر به \(10^{40}\) ارزیابی نیاز دارد.

روش‌های \textbf{\lr{Monte Carlo}} انتگرال‌گیری را با \textbf{نمونه‌گیری}
جایگزین می‌کنند: اگر بتوانیم مقادیر زیادی از پسین بکشیم، هر انتگرالی به
میانگین روی نمونه‌ها تبدیل می‌شود و دقت به تعداد نمونه‌ها بستگی دارد، نه به
تعداد پارامترها. مشکل این است که نمی‌توانیم مستقیماً از پسین نمونه بگیریم؛
معمولاً حتی ثابت نرمال‌سازش را هم نمی‌توانیم حساب کنیم.

\textbf{\lr{Markov chain Monte Carlo (MCMC)}} این مشکل را با ترفندی
به‌طرز زیبایی ساده حل می‌کند: یک گام‌زنی تصادفی در فضای پارامتر که به سمت
مقادیر محتمل‌تر بیشتر از مقادیر کم‌احتمال‌تر حرکت می‌کند، دقیقاً با نسبت‌های
درست. بعد از مدتی، جاهایی که این گام‌زن به آن‌ها سر می‌زند نمونه‌هایی از
پسین‌اند.

\paragraph{این موضوع کجا به کار می‌رود؟}

\begin{itemize}
\tightlist
\item
  هر مدل سلسله‌مراتبی در مقاله‌های اخیر روان‌شناسی و علوم اعصاب با نسخه‌ای
  از آن برازش شده است.
\item
  \textbf{فیزیک و شیمی} (جایی که الگوریتم \lr{Metropolis} در سال 1953
  برای شبیه‌سازی ذرات ابداع شد)، \textbf{همه‌گیرشناسی} (برآورد پارامترهای
  سرایت)، \textbf{تبارشناسی} (بازسازی درخت‌های تکاملی) و \textbf{مالی}
  همه به \lr{MCMC} تکیه دارند.
\item
  \textbf{علوم شناختی نمونه‌گیری را به‌عنوان \emph{نظریه‌ای درباره‌ی ذهن} هم
  به کار می‌برد:} بعضی پژوهشگرها پیشنهاد می‌کنند آدم‌ها استنباط بیزی را با
  نمونه‌گیری ذهنی چند امکان تقریب می‌زنند، که بعضی سوگیری‌ها را توضیح
  می‌دهد.
\item
  \textbf{هوش مصنوعی:} نمونه‌گیری جواب از یک مدل زبانی همان \lr{Monte
  Carlo} است (بخش هوش مصنوعی و \lr{LLM}ها را ببینید).
\end{itemize}

\begin{center}\rule{0.5\linewidth}{0.5pt}\end{center}

\Needspace{11\baselineskip}

\subsection{جزوه‌ی درس}

\paragraph{به زبان ساده}

هر سؤالی که از \lr{posterior} می‌پرسیم («مقدار میانگین چیست؟»، «احتمال
اینکه \(d' > 1\) باشد چقدر است؟») یک انتگرال است، و \lr{grid} نمی‌تواند
انتگرال روی پارامترهای زیاد را حساب کند. \textbf{\lr{Monte Carlo}}
انتگرال را با میانگین جایگزین می‌کند: اگر تعداد زیادی مقدار تصادفی
\emph{کشیده‌شده از} \lr{posterior} داشتیم، میانگین آن‌ها میانگین
\lr{posterior} بود، کسری که بالای 1 است \(P(d' > 1)\) بود، و الی آخر.
مشکل به دست آوردن این نمونه‌هاست. \textbf{الگوریتم \lr{Metropolis}} آن‌ها
را با یک قدم‌زدن تصادفی هوشمندانه به دست می‌آورد. روی یک مقدار پارامتر
بایستید، یک قدم تصادفی کوچک پیشنهاد کنید و مقایسه کنید جای جدید چقدر
معقول است (درست‌نمایی × \lr{prior).} اگر معقول‌تر است، به آنجا بروید. اگر
کمتر معقول است، فقط \emph{گاهی} بروید، با احتمالی برابر نسبت معقول
بودن‌ها؛ وگرنه همان‌جا بمانید و مقدار فعلی را دوباره ثبت کنید. اگر این کار
را به اندازه‌ی کافی ادامه دهید، جاهایی که رفته‌اید نمونه‌ای از
\lr{posterior} هستند. نکته‌ی جالب این است که هرگز به ثابت نرمال‌سازی سخت
نیاز ندارید، چون فقط از نسبت‌ها استفاده می‌شود.

\textbf{زمان‌بندی پیشنهادی (60 دقیقه):} انتگرال‌گیری \lr{Monte Carlo} 12 ·
زنجیره‌های مارکوف 8 · الگوریتم \lr{Metropolis} 18 ·
\lr{Metropolis\lr{--}Hastings} 5 · تنظیم، \lr{burn-in} و خودهمبستگی 12 ·
پیش‌نمایش 5.

\subsubsection{\texorpdfstring{1. انتگرال‌گیری \lr{Monte
Carlo}}{1. انتگرال‌گیری Monte Carlo}}

اگر \(\theta^{(1)}, \dots, \theta^{(S)}\) نمونه‌های مستقلی از
\(p(\theta\mid\mathbf{y})\) باشند، برای هر تابع \(g\):

\[
\mathbb{E}[g(\theta)\mid\mathbf{y}] = \int g(\theta)\, p(\theta\mid\mathbf{y})\, d\theta \;\approx\; \frac{1}{S}\sum_{s=1}^{S} g\big(\theta^{(s)}\big).
\]

\begin{longtable}[]{@{}
  >{\raggedright\arraybackslash}p{(\linewidth - 4\tabcolsep) * \real{0.3333}}
  >{\raggedright\arraybackslash}p{(\linewidth - 4\tabcolsep) * \real{0.3333}}
  >{\raggedright\arraybackslash}p{(\linewidth - 4\tabcolsep) * \real{0.3333}}@{}}
\caption{\centering 1. انتگرال‌گیری \lr{Monte Carlo}}\tabularnewline
\toprule\noalign{}
\begin{minipage}[b]{\linewidth}\raggedright
آنچه می‌خواهیم
\end{minipage} & \begin{minipage}[b]{\linewidth}\raggedright
\(g(\theta)\) را انتخاب کنید
\end{minipage} & \begin{minipage}[b]{\linewidth}\raggedright
برآورد \lr{Monte Carlo}
\end{minipage} \\
\midrule\noalign{}
\endfirsthead
\toprule\noalign{}
\begin{minipage}[b]{\linewidth}\raggedright
آنچه می‌خواهیم
\end{minipage} & \begin{minipage}[b]{\linewidth}\raggedright
\(g(\theta)\) را انتخاب کنید
\end{minipage} & \begin{minipage}[b]{\linewidth}\raggedright
برآورد \lr{Monte Carlo}
\end{minipage} \\
\midrule\noalign{}
\endhead
\bottomrule\noalign{}
\endlastfoot
میانگین پسین & \(\theta\) & میانگین نمونه‌ها \\
\(P(\theta > 1 \mid \mathbf{y})\) & شاخص \(\mathbb{1}[\theta > 1]\) &
نسبت نمونه‌های بالای 1 \\
پسین یک کمیت مشتق‌شده، مثلاً نیمه‌عمر \(2^{1/d} - 1\) & \(g(d)\) & \(g\) را
روی هر نمونه اعمال و بعد خلاصه کنید \\
\lr{credible interval} & \lr{---} & صدک‌ها (یا \lr{HDI)} نمونه‌ها \\
پسین حاشیه‌ای \(\theta_1\) & \lr{---} & فقط ستون‌های دیگر را نادیده بگیرید \\
\end{longtable}

\textbf{خطای استاندارد \lr{Monte Carlo}} میانگین برای نمونه‌های مستقل
\(\mathrm{SD}(g)/\sqrt{S}\) است، مستقل از تعداد پارامترها. به همین دلیل
نمونه‌گیری در ابعاد بالا بر شبکه برتری دارد.

\paragraph{خواندن ریاضی}

\begin{itemize}
\tightlist
\item
  \(\mathbb{E}[g(\theta)\mid\mathbf{y}] = \int g(\theta)\,p(\theta\mid\mathbf{y})\,d\theta\):
  میانگین \lr{posterior} یک کمیت \(g(\theta)\). انتخاب \(g\) یعنی انتخاب
  سؤال (جدول را ببینید).
\item
  \(\approx \frac{1}{S}\sum_{s} g(\theta^{(s)})\): انتگرال را با میانگین
  ساده روی \(S\) نمونه‌ی \(\theta^{(1)}, \dots, \theta^{(S)}\) جایگزین
  کن. بالانویس \((s)\) شماره‌ی نمونه است.
\item
  \textbf{خطا:} \(\mathrm{SD}(g)/\sqrt{S}\). برای نصف کردن خطا چهار
  برابر نمونه لازم است، هر تعداد پارامتری که داشته باشیم.
\item
  \textbf{مثال:} 10,000 نمونه از \(\mathrm{Beta}(8,4)\)؛ 9,000 تای آن‌ها
  بالای 0.5 هستند، پس \(P(\theta > 0.5 \mid \mathbf{y}) \approx 0.90\)
  با خطای \lr{Monte Carlo} برابر
  \(\sqrt{0.9 \times 0.1/10000} = 0.003\).
\end{itemize}

\subsubsection{2. زنجیره‌های مارکوف در یک پاراگراف}

یک \textbf{زنجیره‌ی مارکوف} دنباله‌ای
\(\theta^{(1)}, \theta^{(2)}, \dots\) است که در آن هر مقدار فقط از مقدار
قبلی، با یک قاعده‌ی گذار ثابت، تولید می‌شود. تحت شرایط ملایم، زنجیره نقطه‌ی
شروعش را «فراموش می‌کند» و به یک \textbf{توزیع ایستا} \lr{(stationary)}
می‌رسد: نسبت بلندمدت زمانی که در هر ناحیه می‌گذراند. \lr{MCMC} قاعده‌ی
گذاری می‌سازد که توزیع ایستایش همان پسین است. دو پیامد: - بخش اولیه‌ی
زنجیره هنوز بازتاب نقطه‌ی شروع است. به‌عنوان \textbf{\lr{burn-in}}
\lr{(warm-up)} دور ریخته می‌شود. - نمونه‌های پشت‌سرهم \textbf{همبسته‌اند}،
پس \(S\) نمونه‌ی همبسته اطلاعات کمتری از \(S\) نمونه‌ی مستقل دارند
(اندازه‌ی مؤثر نمونه، هفته‌ی 12).

\subsubsection{\texorpdfstring{3. الگوریتم
\lr{Metropolis}}{3. الگوریتم Metropolis}}

فرض کنید \(\tilde p(\theta) = p(\mathbf{y}\mid\theta)\,p(\theta)\) پسین
\textbf{نرمال‌نشده} (درست‌نمایی × پیشین) باشد. یک \textbf{توزیع پیشنهادی}
متقارن انتخاب کنید، معمولاً یک گام تصادفی نرمال
\(\theta' \sim \mathcal{N}(\theta, \sigma_{\text{\lr{prop}}})\).

\begin{enumerate}
\def\labelenumi{\arabic{enumi}.}
\tightlist
\item
  از یک \(\theta^{(0)}\) شروع کنید.
\item
  برای \(s = 1, 2, \dots, S\):

  \begin{itemize}
  \tightlist
  \item
    پیشنهاد دهید
    \(\theta' \sim \mathcal{N}(\theta^{(s-1)}, \sigma_{\text{\lr{prop}}})\).
  \item
    احتمال پذیرش را حساب کنید: \[
    \alpha_{\text{\lr{acc}}} = \min\left(1,\ \frac{\tilde p(\theta')}{\tilde p(\theta^{(s-1)})}\right).
    \]
  \item
    \(u \sim \mathrm{Uniform}(0, 1)\) بکشید. اگر
    \(u < \alpha_{\text{\lr{acc}}}\)، \textbf{بپذیرید}:
    \(\theta^{(s)} = \theta'\). وگرنه \textbf{رد کنید}:
    \(\theta^{(s)} = \theta^{(s-1)}\) (مقدار فعلی در زنجیره \emph{تکرار}
    می‌شود).
  \end{itemize}
\end{enumerate}

\textbf{چرا به‌طور شهودی کار می‌کند:} حرکت‌های سربالا همیشه پذیرفته می‌شوند؛
حرکت‌های سرپایینی با احتمالی برابر نسبت چگالی‌ها پذیرفته می‌شوند. پس زنجیره
در هر ناحیه به نسبت احتمال پسینش وقت می‌گذراند.

\textbf{چرا فقط پسین نرمال‌نشده لازم است:} شواهد \(p(\mathbf{y})\) هم در
صورت و هم در مخرج نسبت ظاهر می‌شود و ساده می‌شود. این کلید سودمندی
\lr{MCMC} است.

\textbf{جزئیات عملی:} - با \textbf{لگاریتم} کار کنید: اگر
\(\ln u < \ln\tilde p(\theta') - \ln\tilde p(\theta)\) بپذیرید. -
\textbf{پارامترهای کران‌دار:} اگر یک پیشنهاد خارج از تکیه‌گاه بیفتد (مثلاً
\(A > 1\))، پیشینش 0 است، پس \(\ln\tilde p = -\infty\) و پیشنهاد خودکار
رد می‌شود. (راه دیگر نمونه‌گیری روی یک مقیاس تبدیل‌شده و بی‌کران است، مثل
هفته‌ی 3؛ در این صورت باید \lr{Jacobian} را به لگاریتم پسین اضافه کرد.) -
\textbf{چند پارامتر:} همه‌ی پارامترها را هم‌زمان با برداری از انحراف
معیارهای پیشنهادی پیشنهاد دهید، یا هر بار یک پارامتر را به‌روز کنید (که
به نمونه‌گیری \lr{Gibbs} می‌رسد، هفته‌ی 12).

\begin{quote}
\textbf{پیشرفته: تعادل جزئی \lr{(detailed balance).}} یک قاعده‌ی گذار
\(T(\theta \to \theta')\) توزیع \(\pi\) را ایستا نگه می‌دارد اگر برای
همه‌ی جفت‌ها
\(\pi(\theta)\,T(\theta \to \theta') = \pi(\theta')\,T(\theta' \to \theta)\).
این را برای \lr{Metropolis} با \(\tilde p(\theta') < \tilde p(\theta)\)
بررسی کنید: سمت چپ
\(\pi(\theta)\, q(\theta'\mid\theta)\, \tilde p(\theta')/\tilde p(\theta)\)
و سمت راست \(\pi(\theta')\, q(\theta\mid\theta') \cdot 1\) است. با \(q\)
متقارن و \(\pi \propto \tilde p\) این دو برابرند.
\end{quote}

\paragraph{خواندن ریاضی}

\begin{itemize}
\tightlist
\item
  \(\tilde p(\theta) = p(\mathbf{y}\mid\theta)\,p(\theta)\): درست‌نمایی ×
  \lr{prior}، \emph{بدون} تقسیم بر \lr{evidence.} علامت تیلدا
  (\lr{\textasciitilde}) یعنی «نرمال‌نشده».
\item
  \(\alpha_{\text{\lr{acc}}} = \min(1, \tilde p(\theta')/\tilde p(\theta))\):
  اگر پیشنهاد معقول‌تر باشد، نسبت بزرگ‌تر از 1 است و \(\min\) آن را روی 1
  نگه می‌دارد (همیشه پذیرش). اگر نصف معقول باشد، با احتمال 0.5 پذیرفته
  می‌شود.
\item
  «\(u \sim \mathrm{Uniform}(0,1)\) بکش؛ اگر \(u < \alpha\) بود بپذیر»
  فقط راهی است برای اینکه یک رویداد با احتمال \(\alpha\) اتفاق بیفتد (یک
  سکه‌ی نامتعادل).
\item
  \textbf{مثال حل‌شده، قدم 3:} پیشنهاد 0.95 فقط 0.052 برابر 0.80 معقول
  است؛ \(u = 0.30 > 0.052\)، پس رد می‌شود و 0.80 بار دوم در زنجیره نوشته
  می‌شود.
\end{itemize}

\subsubsection{\texorpdfstring{4. \lr{Metropolis\lr{--}Hastings:} پیشنهادهای
نامتقارن}{4. Metropolis\lr{--}Hastings: پیشنهادهای نامتقارن}}

اگر توزیع پیشنهادی متقارن نباشد،
\(q(\theta'\mid\theta) \ne q(\theta\mid\theta')\) (مثلاً پیشنهاد
\lr{log-normal} برای یک نرخ مثبت)، احتمال پذیرش به یک تصحیح نیاز دارد:

\[
\alpha_{\text{\lr{acc}}} = \min\left(1,\ \frac{\tilde p(\theta')\, q(\theta \mid \theta')}{\tilde p(\theta)\, q(\theta' \mid \theta)}\right).
\]

\lr{Metropolis} حالت خاصی است که جمله‌های \(q\) ساده می‌شوند. نمونه‌گیری
\lr{Gibbs} (هفته‌ی 12) حالت خاص دیگری است که در آن هر پیشنهادی پذیرفته
می‌شود.

\subsubsection{5. تنظیم و تشخیص مشکلات یک نمونه‌گیر}

انحراف معیار پیشنهادی \(\sigma_{\text{\lr{prop}}}\) کارایی را کنترل می‌کند:

\begin{longtable}[]{@{}
  >{\raggedright\arraybackslash}p{(\linewidth - 6\tabcolsep) * \real{0.2500}}
  >{\raggedright\arraybackslash}p{(\linewidth - 6\tabcolsep) * \real{0.2500}}
  >{\raggedright\arraybackslash}p{(\linewidth - 6\tabcolsep) * \real{0.2500}}
  >{\raggedright\arraybackslash}p{(\linewidth - 6\tabcolsep) * \real{0.2500}}@{}}
\caption{\centering 5. تنظیم و تشخیص مشکلات یک نمونه‌گیر}\tabularnewline
\toprule\noalign{}
\begin{minipage}[b]{\linewidth}\raggedright
\(\sigma_{\text{\lr{prop}}}\)
\end{minipage} & \begin{minipage}[b]{\linewidth}\raggedright
نرخ پذیرش
\end{minipage} & \begin{minipage}[b]{\linewidth}\raggedright
رفتار
\end{minipage} & \begin{minipage}[b]{\linewidth}\raggedright
\lr{trace plot} چه شکلی است
\end{minipage} \\
\midrule\noalign{}
\endfirsthead
\toprule\noalign{}
\begin{minipage}[b]{\linewidth}\raggedright
\(\sigma_{\text{\lr{prop}}}\)
\end{minipage} & \begin{minipage}[b]{\linewidth}\raggedright
نرخ پذیرش
\end{minipage} & \begin{minipage}[b]{\linewidth}\raggedright
رفتار
\end{minipage} & \begin{minipage}[b]{\linewidth}\raggedright
\lr{trace plot} چه شکلی است
\end{minipage} \\
\midrule\noalign{}
\endhead
\bottomrule\noalign{}
\endlastfoot
خیلی کوچک & خیلی بالا (\lr{\textgreater} 0.9) & گام‌های ریز؛ کاوش کند؛
خودهمبستگی قوی & خطی که آرام سرگردان است \\
خیلی بزرگ & خیلی پایین (\lr{\textless} 0.1) & بیشتر پیشنهادها در نواحی
کم‌احتمال می‌افتند و رد می‌شوند؛ زنجیره گیر می‌کند & پله‌های صاف با فلات‌های
طولانی \\
تقریباً درست & \ensuremath{\approx} 0.2 تا 0.5 & اختلاط \lr{(mixing)} خوب & یک «کرم پشمالو»
\lr{(fuzzy caterpillar)} \\
\end{longtable}

نتایج نظری کلاسیک نرخ پذیرش بهینه‌ی حدود 0.44 را برای یک پارامتر پیشنهاد
می‌کنند که برای پارامترهای زیاد به حدود 0.23 می‌رسد. این‌ها را هدف بدانید،
نه قانون. \(\sigma_{\text{\lr{prop}}}\) را در یک اجرای مقدماتی تنظیم کنید (نه
در اجرایی که نگه می‌دارید). انحراف معیارهای پیشنهادی خوب هم‌مرتبه با
انحراف معیارهای پسین‌اند، و پسین‌های همبسته به پیشنهادهای همبسته (یا
\lr{reparameterization)} نیاز دارند.

\textbf{\lr{Burn-in.}} بخش اولیه‌ی زنجیره را پیش از رسیدن به ناحیه‌ی
پراحتمال دور بریزید؛ \lr{trace plot} نشان می‌دهد کجا.
\textbf{\lr{Thinning}} (نگه داشتن هر \(m\)-امین نمونه) حافظه را کم می‌کند
ولی اطلاعات را دور می‌ریزد. برای استنباط درست لازم نیست.

\textbf{خودهمبستگی.} خودهمبستگی با تأخیر \(h\) همبستگی بین
\(\theta^{(s)}\) و \(\theta^{(s+h)}\) است. خودهمبستگی بالا یعنی کاوش کند
و نمونه‌های مؤثراً مستقل کمتر.

\textbf{همیشه یک نمونه‌گیر جدید را} روی مسئله‌ای با جواب معلوم
\textbf{اعتبارسنجی کنید} (اینجا، \lr{beta-binomial conjugate} و شبکه‌ی
هفته‌ی 10).

\paragraph{برای یک مهندس}

\begin{itemize}
\tightlist
\item
  \textbf{انتگرال‌گیری \lr{Monte Carlo}} همان کاری است که برای تخمین
  \lr{BER} یا احتمال قطعی \lr{(outage)} با شبیه‌سازی می‌کنید: میانگین یک
  تابع نشانگر روی تعداد زیادی تحقق تصادفی. قانون خطای \(1/\sqrt{S}\)
  همان است، و به همین دلیل نرخ‌های خطای پایین شبیه‌سازی‌های خیلی طولانی
  لازم دارند.
\item
  یک زنجیره‌ی \lr{MCMC} یک \textbf{فرایند مارکوف} است که توزیع ایستای آن
  طوری طراحی شده که \lr{posterior} باشد؛ \lr{burn-in} همان حالت گذرا و
  بخش ایستا همان حالت ماندگار است.
\item
  \textbf{خودهمبستگی} زنجیره دقیقاً همان خودهمبستگی یک فرایند تصادفی است.
  نمونه‌های همبسته اطلاعات کمتری دارند، همان‌طور که نمونه‌های نویز همبسته
  کندتر میانگین‌گیری می‌شوند؛ هفته‌ی 12 این را به «اندازه‌ی مؤثر نمونه»
  تبدیل می‌کند.
\item
  \textbf{تنظیم پیشنهاد \lr{(proposal)}} یک مسئله‌ی اندازه‌ی گام است، مثل
  انتخاب اندازه‌ی گام یک فیلتر تطبیقی: خیلی کوچک کند همگرا می‌شود، خیلی
  بزرگ ناپایدار می‌شود (اینجا: بیشتر قدم‌ها رد می‌شوند).
\item
  تفاوت با شبیه‌سازی معمولی: نمی‌توانیم مستقیم از توزیع هدف نمونه‌ی مستقل
  تولید کنیم، پس فرایندی می‌سازیم که به آن همگرا شود.
\end{itemize}

\paragraph{برای مدرس}

\begin{itemize}
\tightlist
\item
  \textbf{زنجیره را روی تخته اجرا کنید.} منحنی \(\theta^7(1-\theta)^3\)
  را بکشید و چهار قدم مثال حل‌شده را اجرا کنید، در حالی که یک دانشجو
  مقدارهای «\(u\)» را تولید می‌کند (یک اپلیکیشن عدد تصادفی کافی است). قدم
  رد‌شده که 0.80 را \emph{تکرار} می‌کند لحظه‌ی کلیدی است؛ بپرسید چرا باید
  ثبت شود.
\item
  \textbf{تشبیهی که جواب می‌دهد:} گردشگری که شب در یک شهر می‌گردد و
  خیابان‌های روشن را ترجیح می‌دهد. خیابان روشن‌تر: همیشه برو. خیابان
  تاریک‌تر: فقط گاهی برو. بعد از یک پیاده‌روی طولانی، زمانی که در هر
  خیابان گذرانده متناسب با روشنایی آن است.
\item
  \textbf{سه \lr{trace plot} را} (خیلی کوچک، خوب، خیلی بزرگ) از بخش
  \lr{D lab} کنار هم نشان دهید. دانشجوها این‌طوری در پنج دقیقه خواندن
  \lr{trace plot} را یاد می‌گیرند.
\item
  \textbf{سؤال‌های رایج.}

  \begin{itemize}
  \tightlist
  \item
    \emph{«چرا فقط از یک \lr{grid} ریز استفاده نکنیم؟»} با 40 پارامتر
    \lr{grid} ناممکن است؛ هزینه‌ی نمونه‌گیری این‌طور رشد نمی‌کند.
  \item
    \emph{«از کجا بدانیم زنجیره همگرا شده؟»} این موضوع هفته‌ی 12 است (چند
    زنجیره، \(\hat R\)).
  \item
    \emph{«چرا \lr{burn-in} را دور بریزیم؟»} نمونه‌های اولیه هنوز نقطه‌ی
    شروع دلبخواه را نشان می‌دهند، نه \lr{posterior} را.
  \end{itemize}
\end{itemize}

\begin{center}\rule{0.5\linewidth}{0.5pt}\end{center}

\subsection{مثال حل‌شده}

\textbf{مدل:} یادآوری \(y = 7\) از \(n = 10\) با پیشین تخت، پس
\(\tilde p(\theta) = \theta^7(1-\theta)^3\) (بدون ثابت‌ها). پسین دقیق:
\(\mathrm{Beta}(8, 4)\) با میانگین 0.667. پیشنهاد:
\(\theta' = \theta + \mathcal{N}(0,\allowbreak 0.1)\).

\begin{landscape}\small
\begin{longtable}[]{@{}
  >{\raggedright\arraybackslash}p{(\linewidth - 16\tabcolsep) * \real{0.1111}}
  >{\raggedright\arraybackslash}p{(\linewidth - 16\tabcolsep) * \real{0.1111}}
  >{\raggedright\arraybackslash}p{(\linewidth - 16\tabcolsep) * \real{0.1111}}
  >{\raggedright\arraybackslash}p{(\linewidth - 16\tabcolsep) * \real{0.1111}}
  >{\raggedright\arraybackslash}p{(\linewidth - 16\tabcolsep) * \real{0.1111}}
  >{\raggedright\arraybackslash}p{(\linewidth - 16\tabcolsep) * \real{0.1111}}
  >{\raggedright\arraybackslash}p{(\linewidth - 16\tabcolsep) * \real{0.1111}}
  >{\raggedright\arraybackslash}p{(\linewidth - 16\tabcolsep) * \real{0.1111}}
  >{\raggedright\arraybackslash}p{(\linewidth - 16\tabcolsep) * \real{0.1111}}@{}}
\caption{\centering مثال حل‌شده}\tabularnewline
\toprule\noalign{}
\begin{minipage}[b]{\linewidth}\raggedright
گام
\end{minipage} & \begin{minipage}[b]{\linewidth}\raggedright
\(\theta\) فعلی
\end{minipage} & \begin{minipage}[b]{\linewidth}\raggedright
پیشنهاد \(\theta'\)
\end{minipage} & \begin{minipage}[b]{\linewidth}\raggedright
\(\tilde p(\theta)\)
\end{minipage} & \begin{minipage}[b]{\linewidth}\raggedright
\(\tilde p(\theta')\)
\end{minipage} & \begin{minipage}[b]{\linewidth}\raggedright
نسبت
\end{minipage} & \begin{minipage}[b]{\linewidth}\raggedright
\(u\)
\end{minipage} & \begin{minipage}[b]{\linewidth}\raggedright
تصمیم
\end{minipage} & \begin{minipage}[b]{\linewidth}\raggedright
\(\theta\) جدید
\end{minipage} \\
\midrule\noalign{}
\endfirsthead
\toprule\noalign{}
\begin{minipage}[b]{\linewidth}\raggedright
گام
\end{minipage} & \begin{minipage}[b]{\linewidth}\raggedright
\(\theta\) فعلی
\end{minipage} & \begin{minipage}[b]{\linewidth}\raggedright
پیشنهاد \(\theta'\)
\end{minipage} & \begin{minipage}[b]{\linewidth}\raggedright
\(\tilde p(\theta)\)
\end{minipage} & \begin{minipage}[b]{\linewidth}\raggedright
\(\tilde p(\theta')\)
\end{minipage} & \begin{minipage}[b]{\linewidth}\raggedright
نسبت
\end{minipage} & \begin{minipage}[b]{\linewidth}\raggedright
\(u\)
\end{minipage} & \begin{minipage}[b]{\linewidth}\raggedright
تصمیم
\end{minipage} & \begin{minipage}[b]{\linewidth}\raggedright
\(\theta\) جدید
\end{minipage} \\
\midrule\noalign{}
\endhead
\bottomrule\noalign{}
\endlastfoot
1 & 0.50 & 0.58 & 0.000977 & 0.001636 & 1.675 & 0.42 & پذیرش (سربالا) &
0.58 \\
2 & 0.58 & 0.80 & 0.001636 & 0.001678 & 1.026 & 0.71 & پذیرش (سربالا) &
0.80 \\
3 & 0.80 & 0.95 & 0.001678 & 0.000087 & 0.052 & 0.30 & \textbf{رد}:
\(u > 0.052\) & 0.80 \\
4 & 0.80 & 0.66 & 0.001678 & 0.002144 & 1.278 & 0.88 & پذیرش (سربالا) &
0.66 \\
\end{longtable}
\end{landscape}

\textbf{زنجیره تا اینجا:} 0.50، 0.58، 0.80، \textbf{0.80}، 0.66. گام
ردشده مقدار 0.80 را در زنجیره \emph{تکرار} می‌کند. این ضروری است: مقادیر
تکراری راهی است که زنجیره ثبت می‌کند 0.80 از 0.95 محتمل‌تر است. حذف آن‌ها
نمونه را سوگیر می‌کند.

\textbf{تفسیر.} در چهار گام، زنجیره از نقطه‌ی شروع به ناحیه‌ی چگالی پسین
بالا (0.6 تا 0.8) رفت و از پرش به دنباله امتناع کرد. در طول هزاران گام،
نسبت زمانی که نزدیک هر مقدار گذرانده می‌شود به \(\mathrm{Beta}(8, 4)\)
همگرا می‌شود.

\begin{center}\rule{0.5\linewidth}{0.5pt}\end{center}

\Needspace{11\baselineskip}

\subsection{\texorpdfstring{\lr{lab} (2 ساعت)}{lab (2 ساعت)}}

\subsubsection{اهداف}

\begin{enumerate}
\def\labelenumi{\arabic{enumi}.}
\tightlist
\item
  انتگرال‌گیری \lr{Monte Carlo} با نمونه‌های مستقل: دیدن خطای
  \(1/\sqrt{S}\).
\item
  نوشتن یک نمونه‌گیر \lr{Metropolis} و اعتبارسنجی آن روی مدل
  \lr{beta-binomial.}
\item
  نمونه‌گیری از پسین مدل فراموشی توانی، مقایسه با شبکه‌ی هفته‌ی 10، و
  محاسبه‌ی یک کمیت مشتق‌شده.
\item
  تنظیم پیشنهاد: مقایسه‌ی انحراف معیارهای پیشنهادی خیلی کوچک، خوب و خیلی
  بزرگ با \lr{trace plot}، نرخ پذیرش و خودهمبستگی.
\end{enumerate}

\textbf{زمان‌بندی:} بخش \lr{A} 15 دقیقه · بخش \lr{B} 25 دقیقه · بخش
\lr{C} 35 دقیقه · بخش \lr{D} 30 دقیقه · بخش \lr{E} 15 دقیقه.

\subsubsection{\texorpdfstring{بخش \lr{A:} انتگرال‌گیری \lr{Monte Carlo}
(15 دقیقه)}{بخش A: انتگرال‌گیری Monte Carlo (15 دقیقه)}}

ده برابر نمونه‌ی بیشتر حدود \(\sqrt{10} \approx 3.2\) برابر خطای
\lr{Monte Carlo} کمتر می‌دهد.

\subsubsection{\texorpdfstring{بخش \lr{C: Metropolis} برای مدل فراموشی
توانی (35 دقیقه)}{بخش C: Metropolis برای مدل فراموشی توانی (35 دقیقه)}}

همان پیشین‌های هفته‌ی 10: \(A \sim \mathrm{Beta}(8, 2)\)،
\(d \sim \mathrm{Gamma}(2, \text{\lr{scale }} 0.3)\).

\textbf{کمیت‌های مشتق‌شده} مجانی به دست می‌آیند. پسین نیمه‌عمر
\(t_{1/2} = 2^{1/d} - 1\) (هفته‌ی 2) و یادآوری پیش‌بینی‌شده بعد از 60 روز
را حساب کنید:

\subsubsection{\texorpdfstring{بخش \lr{D:} تنظیم توزیع پیشنهادی (30
دقیقه)}{بخش D: تنظیم توزیع پیشنهادی (30 دقیقه)}}

\textbf{به این دقت کنید:} زنجیره‌ی «خیلی کوچک» تقریباً همه‌چیز را می‌پذیرد
ولی می‌خزد \lr{(trace} که آرام سرگردان است و خودهمبستگی‌اش حتی با تأخیر 60
نزدیک 1 است)؛ زنجیره‌ی «خیلی بزرگ» تقریباً همه‌چیز را رد می‌کند و مدت‌های
طولانی ساکن می‌ماند (پله‌های صاف). هر دو با 10 000 تکرار برآوردهای غیرقابل
اعتماد می‌دهند: میانگین \(d\) زنجیره‌ی «خیلی کوچک» (0.473) از میانگین پسین
واقعی (0.417) دور است، و میانگین زنجیره‌ی «خیلی بزرگ» فقط به این دلیل
نزدیک است که تقریباً از مقدار شروعش (0.42) تکان نخورده است. زنجیره‌ی خوب
تنظیم‌شده شبیه یک کرم پشمالوست و خودهمبستگی‌اش سریع کم می‌شود.

\subsubsection{\texorpdfstring{بخش \lr{E:} تفسیر (15
دقیقه)}{بخش E: تفسیر (15 دقیقه)}}

آیا خلاصه‌های پسین \lr{Metropolis} برای \((A, d)\) با شبکه‌ی هفته‌ی 10
توافق دارند؟ این توافق چه چیزی را ثابت می‌کند؟ 2. نیمه‌عمر حافظه‌ی این
شرکت‌کننده را، همراه با عدم قطعیت، در یک جمله چطور گزارش می‌کنید؟ 3. چرا
پیشنهادهای ردشده باید به‌صورت مقادیر تکراری در زنجیره ثبت شوند؟

\subsubsection{تمرین‌ها (به ترتیب سختی)}

\textbf{تمرین 1 \ensuremath{\star}.} با نمونه‌گیر بخش \lr{B}،
\(P(\theta > 0.65\mid\mathbf{y})\) را برای 30/50 برآورد کنید و با مقدار
دقیق بخش \lr{A} مقایسه کنید. برای رسیدن به دقت \ensuremath{\pm}0.01 چند تکرار لازم است؟

\textbf{تمرین 2 \ensuremath{\star}.} نمونه‌گیر بخش \lr{C} را از چهار نقطه‌ی شروع خیلی
متفاوت اجرا کنید، مثلاً \((0.3, 0.05)\)، \((0.99, 2.5)\)، \((0.5, 1.0)\)
و \((0.9, 0.1)\). همه‌ی \lr{trace}ها را با هم رسم کنید. بعد از چند تکرار
از هم قابل تشخیص نیستند؟ (این ایده‌ی پشت \lr{R-hat} است، هفته‌ی 12.)

بعد توزیع پیش‌بینی پسین را از نمونه‌هایتان حساب کنید و نمودار \lr{PPC}
هفته‌ی 10 را بازتولید کنید.

\textbf{تمرین 4 \ensuremath{\star}\ensuremath{\star} \lr{(Metropolis\lr{--}Hastings} با پیشنهاد نامتقارن).} از
\(d\) با یک پیشنهاد گام تصادفی \textbf{\lr{log-normal}}،
\(\ln d' \sim \mathcal{N}(\ln d,\allowbreak 0.1)\)، که هرگز مقادیر منفی پیشنهاد
نمی‌دهد، نمونه بگیرید. تصحیح \lr{Hastings} یعنی
\(q(d\mid d')/q(d'\mid d) = d'/d\) را استخراج و اعمال کنید. نشان دهید
حذف این تصحیح پسین را سوگیر می‌کند.

\textbf{تمرین 5 \ensuremath{\star}\ensuremath{\star}\ensuremath{\star} (پیشرفته: تنظیم تطبیقی).} یک مرحله‌ی \lr{warm-up}
بنویسید که هر 500 تکرار انحراف معیارهای پیشنهادی را تنظیم کند تا نرخ
پذیرش به سمت 0.3 برود (مثلاً اگر نرخ خیلی بالاست \lr{SD} را در 1.2 و اگر
خیلی پایین است در 0.8 ضرب کنید)، و بعد آن‌ها را برای اجرای اصلی ثابت کند.
چرا تطبیق باید پیش از نمونه‌هایی که نگه می‌دارید متوقف شود؟

\subsubsection{\texorpdfstring{\lr{lab check-off}}{lab check-off}}

نشان دهید: (1) اعتبارسنجی نمونه‌گیرتان در برابر پسین دقیق \lr{Beta} (بخش
\lr{B)}، (2) \lr{trace plot}ها و خلاصه برای مدل توانی همراه با مقایسه با
شبکه (بخش \lr{C)}، و (3) نمودار تنظیم با نرخ‌های پذیرش (بخش \lr{D).}

\begin{center}\rule{0.5\linewidth}{0.5pt}\end{center}

\subsection{\texorpdfstring{ارتباط با هوش مصنوعی و
\lr{LLM}ها}{ارتباط با هوش مصنوعی و LLMها}}

\textbf{نمونه‌گیری جواب از یک \lr{LLM} همان \lr{Monte Carlo} است.} چون یک
\lr{LLM} خروجی‌اش را نمونه‌گیری می‌کند (کادر هفته‌ی 1)، پرسیدن چندباره‌ی یک
سؤال جواب‌های مختلفی می‌دهد. ترفند رایجی به نام
\textbf{\lr{self-consistency}} تعداد \(k\) جواب نمونه‌گیری می‌کند و رأی
اکثریت را می‌گیرد \lr{{[}Wang et al.,} 2023, \lr{ICLR{]}.} آیا کمک می‌کند،
و چقدر؟ این یک سؤال \lr{Monte Carlo} است: دقت رأی اکثریت را به‌راحتی
نمی‌توانیم به‌صورت بسته حساب کنیم، ولی می‌توانیم شبیه‌سازی‌اش کنیم و خطای
استاندارد \lr{Monte Carlo} را گزارش کنیم، مثل \lr{lab} این هفته.

حالا رأی اکثریت چه می‌کند؟ چرا؟ جوابتان را با فرمول \(\sqrt{p(1-p)/n}\)
بررسی کنید. - \ensuremath{\star}\ensuremath{\star} ارتباط با شناخت: ایده‌ی «خرد جمعی درون فرد» \lr{(wisdom
of the crowd within)} می‌گوید میانگین گرفتن از حدس‌های تکراری یک نفر دقت
را بالا می‌برد.

\begin{center}\rule{0.5\linewidth}{0.5pt}\end{center}

\subsection{اشتباهات رایج و برداشت‌های نادرست}

\begin{enumerate}
\def\labelenumi{\arabic{enumi}.}
\tightlist
\item
  \textbf{حذف گام‌های ردشده} به‌جای تکرار مقدار فعلی. این کار نمونه را به
  سمت دنباله‌ها سوگیر می‌کند.
\item
  \textbf{استفاده از پسین نرمال‌شده.} لازم نیست و معمولاً ممکن هم نیست؛
  شواهد در نسبت ساده می‌شود.
\item
  \textbf{محاسبه‌ی نسبت بدون لگاریتم.} درست‌نمایی‌ها سرریز می‌کنند. لگاریتم
  چگالی‌ها را مقایسه کنید.
\item
  \(\ln\tilde p = -\infty\)
\item
  \textbf{قضاوت درباره‌ی نمونه‌گیر فقط با نرخ پذیرش.} نرخ 30\% همگرایی را
  تضمین نمی‌کند. به \lr{trace}ها و خودهمبستگی نگاه کنید و چند زنجیره اجرا
  کنید (هفته‌ی 12).
\item
  \textbf{نگه داشتن \lr{burn-in}} یا تنظیم پیشنهاد در طول نمونه‌هایی که
  نگه می‌دارید. هر دو نتایج را سوگیر می‌کنند.
\item
  \textbf{«تکرار بیشتر همیشه درستش می‌کند.»} اگر زنجیره در یک مُد گیر کرده
  یا مقیاس پیشنهاد بد است، تکرار بیشتر همان زنجیره ممکن است در زمان
  معقول کمکی نکند. اول تشخیص دهید.
\item
  \textbf{فراموش کردن تصحیح \lr{Hastings}} با پیشنهادهای نامتقارن.
\end{enumerate}

\begin{center}\rule{0.5\linewidth}{0.5pt}\end{center}

\subsection{خودآزمایی}

\textbf{1.} چرا خطای \lr{Monte Carlo} مستقل از تعداد پارامترهاست، ولی
تقریب شبکه‌ای نه؟ \textbf{پاسخ} برآورد \lr{MC} میانگین \(S\) نمونه است و
خطای استانداردش \(\mathrm{SD}/\sqrt{S}\) است، بدون توجه به بُعد
\(\theta\). یک شبکه باید کل فضای پارامتر را بپوشاند، پس اندازه‌اش با
تعداد پارامترها به‌طور نمایی رشد می‌کند.

\textbf{2.} در الگوریتم \lr{Metropolis} چرا به شواهد \(p(\mathbf{y})\)
نیاز نداریم؟ \textbf{پاسخ} پذیرش فقط به \emph{نسبت} چگالی‌های پسین در
مقدار پیشنهادی و فعلی بستگی دارد؛ \(p(\mathbf{y})\) در هر دو ظاهر می‌شود
و ساده می‌شود.

\textbf{3.} \(\tilde p(\theta) = 0.004\) فعلی،
\(\tilde p(\theta') = 0.001\) پیشنهادی و \(u = 0.18\). پذیرش یا رد؟ اگر
\(u = 0.31\) باشد چه؟ \textbf{پاسخ} نسبت 0.25 است. با
\(u = 0.18 < 0.25\) بپذیرید؛ با \(u = 0.31 > 0.25\) رد کنید و مقدار فعلی
را تکرار کنید.

\textbf{4.} نرخ پذیرش یک زنجیره 0.97 است. احتمالاً چه مشکلی هست و
\lr{trace} آن چه شکلی دارد؟ \textbf{پاسخ} انحراف معیار پیشنهادی خیلی
کوچک است: تقریباً هر گام ریزی پذیرفته می‌شود، زنجیره خیلی آرام کاوش می‌کند
و \lr{trace} مثل خطی است که آرام سرگردان است و خودهمبستگی بالایی دارد.
انحراف معیار پیشنهادی را بزرگ کنید.

\textbf{5.} \lr{burn-in} چیست و طولش را چطور تعیین می‌کنید؟ \textbf{پاسخ}
تکرارهای اولیه، پیش از رسیدن زنجیره به ناحیه‌ی پراحتمال، که هنوز بازتاب
مقدار شروع دلخواه‌اند. طولش را از روی \lr{trace plot}ها (و از چند زنجیره
که از نقاط مختلف شروع شده‌اند) انتخاب کنید و آن تکرارها را دور بریزید.

\textbf{6.} تصحیح \lr{Hastings} کی لازم است؟ یک مثال بزنید.
\textbf{پاسخ} وقتی پیشنهاد نامتقارن است،
\(q(\theta'\mid\theta) \ne q(\theta\mid\theta')\)؛ مثلاً پیشنهاد
\lr{log-normal} برای یک نرخ مثبت، یا پیشنهادی که به موقعیت مقدار فعلی در
یک بازه‌ی کران‌دار بستگی دارد.

\textbf{7.} پسین نیمه‌عمر \(t_{1/2} = 2^{1/d} - 1\) را از نمونه‌های
\lr{MCMC} \(d\) چطور به دست می‌آورید؟ \textbf{پاسخ} هر نمونه‌ی \(d\) را به
یک نیمه‌عمر تبدیل کنید و نمونه‌های تبدیل‌شده را خلاصه کنید (میانگین، صدک‌ها،
\lr{HDI).} هیچ استخراج اضافه‌ای لازم نیست؛ این یکی از مزیت‌های عملی بزرگ
نمونه‌گیری است.

\begin{center}\rule{0.5\linewidth}{0.5pt}\end{center}

\subsection{تکلیف}

این هفته تکلیف نمره‌دار جدیدی نداریم. \textbf{\lr{A4} (تحلیل بیزی با
شبکه)} در ابتدای درس هفته‌ی 12 تحویل داده می‌شود.

\textbf{تمرین اختیاری (بدون نمره):} تمرین 2 (چند نقطه‌ی شروع) بهترین
آمادگی برای هفته‌ی 12 است.

\begin{center}\rule{0.5\linewidth}{0.5pt}\end{center}

\subsection{مطالعه‌ی بیشتر}

\begin{itemize}
\tightlist
\item
  \textbf{\lr{Kruschke} (2014)، فصل 7.} \lr{Markov chain Monte Carlo:}
  مقدمه‌ای شهودی بر الگوریتم \lr{Metropolis} (مثال «سیاستمدار»)، منطق آن
  و نماینده بودن نمونه‌ها. \emph{(منبع اصلی.)}
\item
  \textbf{\lr{Farrell} \lr{\&} \lr{Lewandowsky} (2018)، فصل 7.} برآورد بیزی
  پارامتر با \lr{MCMC}، به‌کاررفته روی مدل‌های شناختی.
\item
  \textbf{\lr{Martin, Kumar} \lr{\&} \lr{Lao} (2021)}: بخش‌های موتورهای
  استنباط و \lr{MCMC} {[}شماره‌ی فصل بررسی شود{]}.
\end{itemize}

\emph{پیش‌نمایش هفته‌ی 12:} نمونه‌گیری \lr{Gibbs} (بدون هیچ ردی)، و سؤالی
که هر کاربر \lr{MCMC} باید جواب دهد: آیا زنجیره‌ام همگرا شده است؟
